% ===================================================================
%  TABELO N° 9 — La monto. La foresto.
%  Feuillets 63 a 69 du fac-simile = folios 61 a 67.
%  Correspondance : folio imprime = numero de feuillet - 2.
%
%  Ca tabelo esas en du ceni, kada un kun sua propra intitulo di
%  personi e sua propra numerado d'alineas, ripetanta de 1 : «Unesma
%  ceno» (La Monto. — La Balnurbo) e «Duesma ceno» (La Foresto. — La
%  Chasado). La cles %%K sequas la numerado quale imprimita.
% ===================================================================

% --- Feuillet 63 (folio 61) -----------------------------------------
\begin{VUpage}[63]{61}
%%K t09-apar-1 apar
\VUsaut{4.0mm}
\VUtitre{15.0pt}{60}{TABELO N\textsuperscript{o} 9}
\VUsaut{3.0mm}

%%K t09-tit-1 sub
\VUcentre{12.0pt}{0}{\textit{Unesma ceno.}}
\VUsaut{3.0mm}

%%K t09-tit-2 sub
\VUpk{11.4pt}{40}{La Monto. --- La Balnurbo.}
\VUsaut{3.0mm}

%%K t09-c1-01-1 p
\VUblancAlinea
1. --- De ok dii me esas en Pratz. Me ne\nl
povas expresar la tota admiro quan me sentis\nl
kande me esis avan la marveloza \VUgras{montaro} \textsuperscript{(1)}\nl
dil Pirenei di qua la \VUgras{kresto} \textsuperscript{(2)} aspektas sur la\nl
blua cielo quale giganta festono.

%%K t09-c1-02-1 p
\VUblancAlinea
2. --- Me ne imaginis, ke la Pireneal \VUgras{so}\cc
\VUgras{miti} \textsuperscript{(3)} atingas \VUgras{alteso} tante granda, e me restis\nl
astonegata koram la belega \VUgras{glacieri} \textsuperscript{(5)} e la\nl
\VUgras{monto-pinti} \textsuperscript{(4)} nivoza sur qui rulas tante tero\cc
rigant \VUgras{avalanchi}.

%%K t09-tit-3 sub
\VUsaut{3.0mm}
\VUpk{10.2pt}{40}{Monto-peizajo.}
\VUsaut{3.0mm}

%%K t09-c1-03-1 p
\VUblancAlinea
3. --- Quik de la morga dio di mea arivo, me\nl
iris a la olda \VUgras{volkano} \textsuperscript{(6)} extingita qua esas\nl
proxim Pratz. Sur la bordi arida e nudigita dil\nl
\VUgras{kratero}, on vidas ankore bone la traci dil paso\nl
di la \VUgras{lavao} qua fluis, plura yarcenti ante nun,\nl
sur la \VUgras{flanko dil monto} \textsuperscript{(7)}. Me sucesis trovar,\nl
sur un de la \VUgras{penti}, kelka fragmenti di ta lavao.

%%K t09-c1-04-1 p
\VUblancAlinea
4. --- Pratz esas situita sur la frontiero, e la\nl
\VUgras{fortreso} \textsuperscript{(8)} qua defensas la \VUgras{defileo} \textsuperscript{(27)} separanta\nl
Francia de Hispania tote memorigas ta olda\nl
kasteli di la rivo di Rheno qui semblas guatar\nl
\VUgras{raptajo} de la supro di sua roko.

%%K t09-c1-05-1 p
\VUblancAlinea
5. --- Enorm \VUgras{aglo} \textsuperscript{(9)} qua havas sua nesto ne\nl
fore de la \VUgras{fuorto} \textsuperscript{(8)} ofte atraktas mea atenco.
\end{VUpage}

% --- Feuillet 64 (folio 62) -----------------------------------------
\begin{VUpage}[64]{62}
\VUblancAlinea
Ta \VUgras{raptera ucelo}, di qua la \VUgras{beko} esas forta,\nl
ofte adportas a sua yuni mutonyuno o kaproyu\cc
no quin lu tenas per sua \VUgras{ungli.} Tanta esas la\nl
grandeso di lua \VUgras{ali} ke lu povas longatempe\nl
glitflugar kun lua pezoza portajo.

%%K t09-tit-4 sub
\VUsaut{3.0mm}
\VUpk{10.2pt}{40}{L'Exkursionero.}
\VUsaut{3.0mm}

%%K t09-c1-06-1 p
\VUblancAlinea
6. --- Ibe la maxim agreabla promeno esas\nl
l'exkurso a la kaskado \textsuperscript{(10)} di « Kau ». L'\VUgras{aquo}\cc
\VUgras{falo} vortice falas e pose desaparas kom bela\nl
\VUgras{rivereto} en \VUgras{bosko de abieti} \textsuperscript{(11)} situita weste\nl
de la urbo. Ni acensis tra ica bosko til la \VUgras{me}\cc
\VUgras{teorologial observatorio} \textsuperscript{(12)} e de la supro dil\nl
\VUgras{belvedero}, ni parvidis la tota \VUgras{panoramo} di la\nl
regiono. La \VUgras{peizajo} esas marveloza, e la \VUgras{na}\cc
\VUgras{turo} semblas prodigir a la lando omna trezori\nl
quin lu disponas.

%%K t09-c1-07-1 p
\VUblancAlinea
7. --- Por decensar, ni uzis la \VUgras{funikularo} \textsuperscript{(13)}\nl
e ni haltis en mikra \VUgras{staciono} proxim la \VUgras{rui}\cc
\VUgras{naji} \textsuperscript{(14)} di olda \VUgras{feudala kastelo} olim habitita\nl
dal siniori di Pratz.

%%K t09-c1-08-1 p
\VUblancAlinea
8. --- Kompatinda kastelo ! Quon kredeble ol\nl
pensas vidante infre la koketa \VUgras{baln-urbo} \textsuperscript{(15)}\nl
qua nun esas \VUgras{termala staciono} segunmoda ?

%%K t09-c1-08-2 p
\VUblancAlinea
La \VUgras{klimato} esas tante agreabla, la \VUgras{mineral}\nl
\VUgras{aquo} tante efikanta, ke la \VUgras{kuracado} uzata en la\nl
\VUgras{sanatorio} \textsuperscript{(17)} esas vera plezuro.

%%K t09-tit-5 sub
\VUsaut{3.0mm}
\VUpk{10.2pt}{40}{Agreabla termal urbo.}
\VUsaut{3.0mm}

%%K t09-c1-09-1 p
\VUblancAlinea
9. --- Cetere, on agis omno por igar agreabla\nl
la sejorno. Granda \VUgras{viadukto} \textsuperscript{(16)} esis konstruk\cc
tita por posibligar fervoyo; la \VUgras{parko} \textsuperscript{(18)} di la
\parplein
\end{VUpage}

% --- Feuillet 65 (folio 63) -----------------------------------------
\begin{VUpage}[65]{63}
%%K t09-c1-09-1 p suite
\VUcontinue
\VUgras{termal establisuro}, en qua on \VUgras{koncertas}, esas\nl
aranjita en charmanta maniero. Plura gracioz\nl
« \VUgras{villa} »-i \textsuperscript{(19)} esis konstruktita cirkum la ka\cc
zino \textsuperscript{(21)}, e \VUgras{choseo} \textsuperscript{(22)} tre bone mantenata\nl
grande faciligas l'interkomuniki.

%%K t09-c1-10-1 p
\VUblancAlinea
10. --- Simpla \VUgras{taluso} \textsuperscript{(23)} separas la \VUgras{voyo} \textsuperscript{(20)}\nl
de la \VUgras{valo} \textsuperscript{(25)} propre dicita, e funde di qua\nl
jacas gracioza mikra \VUgras{lago} \textsuperscript{(26)} qua alimentas\nl
\VUgras{rivereto} \textsuperscript{(24)} limpida.

%%K t09-c1-11-1 p
\VUblancAlinea
11. --- Sur l'altra latero dil valo, \VUgras{fer-mi}\cc
\VUgras{neyo} \textsuperscript{(28)} esas explotata. Plurfoye me iris vidar\nl
la \VUgras{ministi} decensanta en la \VUgras{shakto}, e mem me\nl
eniris la \VUgras{galerio} en qua li pasas sua jorno, ex\cc
traktante la \VUgras{erco} qua kontenas la \VUgras{metalo}.

%%K t09-c1-12-1 p
\VUblancAlinea
12. --- Importanta \VUgras{fuzerio} esis konstruktita\nl
proxim la mineyo por quik utiligar la richaji\nl
qui esas extraktata de olu. La \VUgras{fornego} \textsuperscript{(29)}, var\cc
migita per la \VUgras{ter-karbono} qua abundas hike,\nl
posibligas rapide e chipe obtenar \VUgras{gisfero}.

%%K t09-c1-13-1 p
\VUblancAlinea
13. --- Ne fore de la mineyo on exkavas\nl
\VUgras{petro-mineyo} \textsuperscript{(30)}. Nek \VUgras{graniton}, nek \VUgras{porfiron},\nl
nek mem \VUgras{greson} la olda \VUgras{petro-ministo} detran\cc
chas per sua \VUgras{pikono}, ma simpla \VUgras{quadron} por\nl
la konstrukto di domi.

%%K t09-c1-14-1 p
\VUblancAlinea
14. --- Un de la maxim bela ornivi di ta\nl
lando esas la \VUgras{abieti} \textsuperscript{(31)}. Omnaloke en la fundo\nl
di \VUgras{ravino} \textsuperscript{(32)}, sur la rivo di \VUgras{torento} \textsuperscript{(33)}, di\nl
\VUgras{abismo} \textsuperscript{(34)} o di precipiso, lia dina aguli verde\nl
nuancizas.

%%K t09-tit-6 sub
\VUsaut{3.0mm}
\VUpk{10.2pt}{40}{Pedirado.}
\VUsaut{3.0mm}

%%K t09-c1-15-1 p
\VUblancAlinea
15. --- Me profitas mea vakanci por \VUgras{longe}\nl
promenar. La \VUgras{voyi} \textsuperscript{(35)} qui sinuifas sur la monto
\parplein
\end{VUpage}

% --- Feuillet 66 (folio 64) -----------------------------------------
\begin{VUpage}[66]{64}
%%K t09-c1-15-1 p suite
\VUcontinue
posibligas parkurar, sen egardar la disto, plura\nl
\VUgras{kilometri} ed ofte plura \VUgras{quaropa kilometri}.

%%K t09-c1-15-2 p
\VUblancAlinea
\VUgras{Termini} \textsuperscript{(36)} ed \VUgras{indik-fosti} \textsuperscript{(37)} markizas la\nl
voyi sur qui ofte on renkontras yuna \VUgras{monta}\cc
\VUgras{no} \textsuperscript{(38)} kun \VUgras{bisako} \textsuperscript{(40)} an lua flanko, portante\nl
\VUgras{marmoto} \textsuperscript{(39)}, o kelka \VUgras{cigano} \textsuperscript{(41)} di qua la \VUgras{fur}\cc
\VUgras{toko} \textsuperscript{(42)} stranje kontrastas kun la olda \VUgras{ka}\cc
\VUgras{puco} \textsuperscript{(43)} lacerita. Il duktas per \VUgras{kateno urso} \textsuperscript{(44)}\nl
dresita.

%%K t09-tit-7 sub
\VUpk{10.2pt}{40}{Monti-acenso.}
\VUsaut{3.0mm}

%%K t09-c1-16-1 p
\VUblancAlinea
16. --- Preske singladie me iras kun \VUgras{guidis}\cc
\VUgras{to} \textsuperscript{(48)} praktikar \VUgras{acenso}, e me esas tre fiera\nl
kande, kun \VUgras{sako} \textsuperscript{(47)} an la dorso, e la « \VUgras{alpens}\cc
\VUgras{tock} » en manuo, quale vera \VUgras{turisto} \textsuperscript{(46)} me\nl
asaltas la \VUgras{rokaji} \textsuperscript{(45)}.

%%K t09-c1-17-1 p
\VUblancAlinea
17. --- De la supro di monto, la habitanti dil\nl
planajo aspektas ne plu grosa kam formiki;\nl
\VUgras{muton-trupo} \textsuperscript{(50)} pasturanta en \VUgras{pastureyo} \textsuperscript{(49)}\nl
semblas ludilo por infanti. \VUgras{Pino} \textsuperscript{(51)} od anke\nl
\VUgras{ligna dometo} \textsuperscript{(52)} aspektas apene kom makulo\nl
sur la verdajo.

%%K t09-c1-18-1 p
\VUblancAlinea
18. --- Dum ta exkursi, ni ofte renkontras\nl
sur la \VUgras{voyeto} \textsuperscript{(53)} \VUgras{chamo-chasisto} \textsuperscript{(54)} shultre\nl
portanta la bestio quan il jus mortigis.

%%K t09-c1-19-1 p
\VUblancAlinea
19. --- Me lokizis en mea herbario tre remar\cc
kinda brancho de \VUgras{filiko} \textsuperscript{(55)} e bela rozea bran\cc
cheto di \VUgras{eriko} \textsuperscript{(57)} quin me koliis ye l'enireyo\nl
di mikra \VUgras{groto} \textsuperscript{(56)} lor mea lasta promeno.

%%K t09-tit-8 sub
\VUsaut{3.0mm}
\VUcentre{12.0pt}{0}{\textit{Duesma ceno.}}
\VUsaut{3.0mm}

%%K t09-tit-9 sub
\VUpk{11.4pt}{40}{La Foresto. --- La Chasado.}
\VUsaut{3.0mm}

%%K t09-c2-01-1 p
\VUblancAlinea
1. --- Hiere me pasis jorno delicoza en la fo\cc
resto. La diligenteso qua regnas inter la ani\ccplein
\end{VUpage}

% --- Feuillet 67 (folio 65) -----------------------------------------
\begin{VUpage}[67]{65}
%%K t09-c2-01-1 p suite
\VUcontinue
mali e l'\VUgras{insekti} \textsuperscript{(5)} da qui ol habitesas tre inte\cc
resis me.

%%K t09-c2-01-2 p
\VUblancAlinea
L'\VUgras{araneo} \textsuperscript{(1)} qua texas sua \VUgras{reto} \textsuperscript{(2)} inter du\nl
branchi por kaptar la \VUgras{musho} \textsuperscript{(3)} trairanta, la\nl
\VUgras{heliko} \textsuperscript{(4)} penipenoze portanta sua konko, la\nl
lukano exploranta pri sua kaptajo, la diligenta\nl
formiko tre zeloza proxim la \VUgras{formikeyo} \textsuperscript{(6)}, ta\nl
omna mikri iris, venis, laboris kun extraordi\cc
nara diligenteso.

%%K t09-c2-02-1 p
\VUblancAlinea
2. --- \VUgras{Serpento} \textsuperscript{(7)} quan me unesmavide opi\cc
nionis \VUgras{vipero}, ma qua ne esis altro kam simpla\nl
\VUgras{kolubro}, esis blotisinta sub arboro-trunko, e\nl
tempope ekirigis sua dina kapeto qua sempre\nl
semblis pronta por mordar. Avan me grosa\nl
\VUgras{limako} \textsuperscript{(8)} pezoze reptis pos devorir un de la\nl
saporoza \VUgras{frageti} \textsuperscript{(11)} qui abunde kreskas ibe.

%%K t09-c2-03-1 p
\VUblancAlinea
3. --- La sulo esis tapetizita per \VUgras{boskal flori};\nl
\VUgras{primuli} \textsuperscript{(9)}, \VUgras{pervinki} \textsuperscript{(10)}, e mult altra planti\nl
quin me ne konocis, florifis inter \VUgras{herbo}\cc
\VUgras{tufi} \textsuperscript{(12)}.

%%K t09-c2-04-1 p
\VUblancAlinea
4. --- En ta regiono, la \VUgras{truflo} esas preske\nl
tam komuna kam la \VUgras{fungo} \textsuperscript{(14)} e \VUgras{porkyuno} \textsuperscript{(13)}\nl
apartenanta a forestisto, gurmande exploris\nl
kad lu sucesos trovar kelki de oli.

%%K t09-tit-10 sub
\VUsaut{3.0mm}
\VUpk{10.2pt}{40}{Furtechasado.}
\VUsaut{3.0mm}

%%K t09-c2-05-1 p
\VUblancAlinea
5. --- Me asistis la \VUgras{fino} di ceno qua tre amu\cc
zis me. Per helpo di la flarado di sua \VUgras{hundo ba}\cc
\VUgras{sota} \textsuperscript{(15)}, la \VUgras{forestisto} \textsuperscript{(16)} jus surprizis \VUgras{furte}\cc
\VUgras{chasero} \textsuperscript{(22)}, ye l'instanto kande ica preparis\nl
su kunportar la \VUgras{vildo} \textsuperscript{(17)} quan il mortigabis.\nl
Preferante abandonar sua kaptajo kam lasar
\parplein
\end{VUpage}

% --- Feuillet 68 (folio 66) -----------------------------------------
\begin{VUpage}[68]{66}
%%K t09-c2-05-1 p suite
\VUcontinue
su arestesor, la furte-chasero fugabis, e \VUgras{le}\cc
\VUgras{poro} \textsuperscript{(18)}, \VUgras{cervino} \textsuperscript{(19)}, \VUgras{kapreolo} \textsuperscript{(20)} ed \VUgras{apro} \textsuperscript{(21)}\nl
jacis sur la herbo joyigante la forestisto.

%%K t09-c2-06-1 p
\VUblancAlinea
6. --- La diverseso dil arbori quin kontenas\nl
la foresto esas astonanta. La \VUgras{fagi} \textsuperscript{(23)} e la \VUgras{bir}\cc
\VUgras{ki} \textsuperscript{(26)}, la \VUgras{pini}, qui produktas la \VUgras{rezino} \textsuperscript{(27)}, la\nl
\VUgras{querki} \textsuperscript{(29)} ed anke mult altri intermixas sua\nl
brancharo.

%%K t09-c2-06-2 p
\VUblancAlinea
La \VUgras{hedero} \textsuperscript{(24)}, qua adheras al trunko di la\nl
alt arbori, kolorizas la \VUgras{silvo} \textsuperscript{(25)} per verdajo\nl
brilanta, qua agreable unionas su al nuanco\nl
plu klara e pale verda dil \VUgras{musko} \textsuperscript{(31)} en qua la\nl
\VUgras{pikverdo} \textsuperscript{(30)} serchas insekti.

%%K t09-tit-11 sub
\VUsaut{3.0mm}
\VUpk{10.2pt}{40}{Forestal peizajo.}
\VUsaut{3.0mm}

%%K t09-c2-07-1 p
\VUblancAlinea
7. --- La olda querki charmis me. Lia grosa\nl
\VUgras{radiki} \textsuperscript{(32)} tordita, mi-ekirinta del sulo, donas\nl
a li splendid aspekto di forteso e vigoro, e me\nl
questionas me quale la \VUgras{proprietero} \textsuperscript{(35)} povas\nl
rezignar pri abatigar e livrar oli a la \VUgras{segilo} \textsuperscript{(34)}\nl
di la \VUgras{hakisto} \textsuperscript{(33)}. La \VUgras{trunki} \textsuperscript{(36)} kompatinda, spo\cc
liita de sua \VUgras{kortico} \textsuperscript{(37)} o fendita da la \VUgras{hakilo} \textsuperscript{(38)}\nl
dil \VUgras{jorn-laboristo} \textsuperscript{(39)} chagrenigas me.

%%K t09-c2-08-1 p
\VUblancAlinea
8. --- Proxime, olda \VUgras{karbonifisto} \textsuperscript{(41)} preparis\nl
per sua \VUgras{shovelo amaso} \textsuperscript{(42)} de karbono, e oldino\nl
kunportis \VUgras{fasko} \textsuperscript{(40)} de \VUgras{ligno mortinta} ek bran\cc
cheti dil abatita arbori.

%%K t09-tit-12 sub
\VUsaut{3.0mm}
\VUpk{10.2pt}{40}{Chasado.}
\VUsaut{3.0mm}

%%K t09-c2-09-1 p
\VUblancAlinea
9. --- Me volis irar por repozo dum kelk ins\cc
tanti en la \VUgras{forestistal domo} \textsuperscript{(46)} ma, kande me\nl
arivis proxim la \VUgras{lageto} \textsuperscript{(43)}, sur qua natas bela
\parplein
\end{VUpage}

% --- Feuillet 69 (folio 67) -----------------------------------------
\begin{VUpage}[69]{67}
%%K t09-c2-09-1 p suite
\VUcontinue
\VUgras{cigno} \textsuperscript{(45)}, me remarkis, ke recent \VUgras{inundo} \textsuperscript{(44)}\nl
chanjabis la voyo a vera \VUgras{marsho}. Me mustis\nl
deturnar me, e, pasante proxim la \VUgras{cedro} \textsuperscript{(49)},\nl
la \VUgras{popli} \textsuperscript{(48)}, e la \VUgras{ulmo} \textsuperscript{(47)} qui esas an la bordo\nl
di la foresto, me vidis debushar ek \VUgras{kopso} \textsuperscript{(54)}\nl
belega \VUgras{cervulo} \textsuperscript{(50)}, persequata da obstinanta\nl
\VUgras{(chas)-hundaro} \textsuperscript{(51)} quan ecitis anke la \VUgras{kor}\cc
\VUgras{no} \textsuperscript{(53)} dil \VUgras{chaseri kavalkanta} \textsuperscript{(52)}. La kompa\cc
tinda bestio esis tote extenuita. Ol venis mor\cc
tante falar ye kelka pazi proxim me.

%%K t09-c2-10-1 p
\VUblancAlinea
10. --- La filiulo dil forestisto, quan la brui\cc
sado dil \VUgras{kuro-chaso} atraktabis, ekiris la hemo\nl
e ni due retrovenis aden la foresto. Il vidabis\nl
\VUgras{pigo} \textsuperscript{(56)} sur brancho di olda querko. Il opi\cc
nionis, ke olua nesto kredeble esas celata en\nl
la \VUgras{foliaro} \textsuperscript{(58)} ed il volis kaptar la nesto.

%%K t09-c2-10-2 p
\VUblancAlinea
Ajila quale \VUgras{skurelo} \textsuperscript{(61)}, il klimis de \VUgras{bran}\cc
\VUgras{cho} \textsuperscript{(55)} a brancho, ma il trovis nulo e sucesis\nl
nur igar flugeskar olda \VUgras{korniko} \textsuperscript{(60)} perchinta\nl
sur \VUgras{ramo} \textsuperscript{(57)}.

\VUsaut{4.0mm}
\VUfilet{22mm}
\end{VUpage}
